\documentclass[12pt, twoside]{report}

\usepackage{multicol}
%% comment this line for handouts. Uncomment for lecture notes
%\newcommand{\lectureNoteMode}{}

\raggedbottom

\usepackage{color, xifthen, amsthm, amssymb, amsmath, tikz, pgfplots}
\pgfplotsset{compat=1.18}
\usetikzlibrary{backgrounds}
\usepackage[letterpaper, margin=1in]{geometry}
\usepackage{multicol}
\usepackage{fourier}  %don't use txfonts b/c the sf doesn't look good

\usepackage{enumitem}
\setlist{topsep=-\parskip/2}


\setlength{\parindent}{0pt}
\setlength{\parskip}{\baselineskip}


\usepackage{mdframed} % makes better frames around examples (eg pagebreak allowed)
  % from mdframed package. used for examples
\newmdenv[
  skipabove=10pt,%\topsep,
  skipbelow=10pt,%\topsep
  ]{allrules}

\usepackage{tcolorbox} % for colored boxes


\usepackage{subfigure}
\usepackage{hyperref}
\hypersetup{colorlinks, urlcolor=red, linkcolor=.} %color urls, but not interal refs
\usepackage{graphicx, floatflt}
\graphicspath{{./figs/}{../figs/}}
\DeclareGraphicsExtensions{.png,.jpg,.pdf}

\usepackage{wrapfig}  %useful for wrapping text around pictures or tables
		      %though it doesn't do a great job inside my example boxes


\definecolor{defbod}{rgb}{.8,.8,.8}     \definecolor{defhead}{rgb}{.7,.1,.1}
\definecolor{thmbod}{rgb}{.8,.8,.8}     \definecolor{thmhead}{rgb}{.1,.1,.9}
\definecolor{factbod}{rgb}{1,1,1}    \definecolor{facthead}{rgb}{.1,.7,.1}
\definecolor{exbod}{rgb}{1,1,1}         \definecolor{exhead}{rgb}{0,.3,0}
\definecolor{soltext}{rgb}{.2,.2,.5}
\definecolor{notebod}{rgb}{.9,.9,.9}    \definecolor{notehead}{rgb}{.1,.1,.7}

\newcommand{\defn}[1]{
 \noindent \colorbox{defbod}{\begin{minipage}{.95\linewidth}
 \hspace{-12pt}\colorbox{defhead}{\sf\color{white} Definition}
 #1
 \end{minipage} } \medskip
 }
\newcommand{\dt}[1]{\textcolor{defhead}{\bfseries #1}}

\newcommand{\thrm}[1]{
 \noindent \colorbox{thmbod}{\begin{minipage}{.95\linewidth}
 \hspace{-12pt}\colorbox{thmhead}{\sf\color{white}Theorem}
 #1
 \end{minipage} } \medskip
 }
\newcommand{\coro}[1]{
 \noindent \colorbox{thmbod}{\begin{minipage}{.95\linewidth}
 \hspace{-12pt}\colorbox{thmhead}{\sf\color{white}Corollary}
 #1
 \end{minipage} } \medskip
 }
\newcommand{\fact}[1]{  %%%% !! note this different from calcnotes
 \noindent \begin{allrules}
 #1
 \end{allrules} \medskip
 }

\newcommand{\note}[1]{
	\noindent \colorbox{notebod}{\begin{minipage}{.95\textwidth}
			\hspace{-12pt}\colorbox{notehead}{\sf\color{white}NOTE}
			#1
	\end{minipage} } \medskip
}


% making the example thing with a line down the left.
\newmdenv[
rightline=false,
bottomline=false,
skipabove=0pt,%\topsep,
skipbelow=0pt,%\topsep
]{siderules}

\newcommand{\ex}[1]{
  \noindent \begin{siderules}
    \textcolor{exhead}{\sf Example}
    #1
  \end{siderules} 
  \medskip
}




%%%%%%%%%% set it up so that these are not displayed when 
%%%%%%%%%% printing notes for students?  or for overhead?
\newcommand{\exsol}[1]{
  {\color{soltext} #1}
}

\newcommand{\mynonumsect}[2][]{
 \cleardoublepage
 \section*{#1\hskip16pt #2}
 \addcontentsline{toc}{section}{#1 #2}
}

\newcommand{\myhead}[1]{
\par%\vskip\baselineskip
\noindent
{\large\bf #1}
}

%%%%%%%%%%%%%%%%%%%%%%%%
% symbol shortcuts

\newcommand{\nice}{\displaystyle}
\newcommand{\eqq}{\ensuremath \stackrel{?}{=}}
\newcommand{\R}{\ensuremath{\mathbb{R}}}




% %%%%%%%%%%%%%%%%%%%%%%%%
% %
% % Creating handouts vs notes
% %% If you use the line below before loading this file then it will be set up for
% %% lecture notes.  Default is for handouts
% %%
% %% \newcommand{\lectureNoteMode}{}
% %%
% %% do not uncomment here, just add the line to the file where you input this one.
% \ifthenelse{\not\isundefined\lectureNoteMode}
% { % lecture notes 
%   \newcommand{\Hvskip}[1]{}
%   \newcommand{\Hvfill}{}
%   \newcommand{\Hpbreak}{}  % pagebreaks for handouts
%   \newcommand{\Npbreak}{\vfill\pagebreak} % pagebreaks for printable notes
%   \newcommand{\Nonly}[1]{#1} % shows in notes, not in handouts
%   \newcommand{\mycleardoublepage}{\cleardoublepage}

%   \newcommand{\ex}[1]{
%    \noindent \begin{allrules}
%     \textcolor{exhead}{\sf Example}\par
%     #1
%     \end{allrules} \medskip
%   }

% } 
% { 
% %% Handout mode 
% %% this section sets it up with smaller margins, and blank spots for writing solns to examples
%   %%%%%%%%%%%%%%%%%%%%%%%%%%%%%%%%%%%%%%%%%%
%   %change frame around examples
%   % from mdframed package
%   \newmdenv[
%   rightline=false,
%   bottomline=false,
%   skipabove=0pt,%\topsep,
%   skipbelow=0pt,%\topsep
%   ]{siderules}
  
%   \newcommand{\ex}[1]{
%     \noindent \begin{siderules}
%       \textcolor{exhead}{\sf Example}
%       #1
%     \end{siderules} 
%     \medskip
%   }
  
%   %
%   %%%%%%%%%%%%%%%%%%%%%%%%%%%%%%%%%%%%%%%%%%
  
%   % don't show solutions for examples
%   \renewcommand{\exsol}[1]{}
  
%   %used for making blank space in handouts. eg: \vs2
%   \newcommand{\Hvskip}[1]{\vskip#1in\null}
%   \newcommand{\Hvfill}{\vfill} %vfill only in handouts mode
%   \newcommand{\Hpbreak}{\pagebreak} % pagebreaks for handouts
%   \newcommand{\Npbreak}{} % pagebreaks for printable notes
%   \newcommand{\Nonly}[1]{} % shows in printable, not in handouts
%   \newcommand{\mycleardoublepage}{\newpage} % don't need blank pages if not printing
  
%   % don't really need margins since not printing.
%   \geometry{margin=0.75in}

%   \renewcommand{\mynonumsect}[2][]{
%     \mycleardoublepage
%     \section*{#1\hskip16pt Handout: #2}
%     \addcontentsline{toc}{section}{#1 #2}
%   }

% } 





\setlength{\parindent}{0pt}

\setcounter{chapter}{1}
\begin{document}


\mynonumsect[1.1]{Functions and Notation}

\textbf{Goal:} To evaluate function values

\myhead{What is a Function?}

The natural world is full of relationships between quantities that change.  When we see these relationships, it is natural for us to ask "If I know one quantity, can I then determine the other?" This establishes the idea of an input quantity, or \dt{independent variable}, and a corresponding output quantity, or \dt{dependent variable}. From this we get the notion of a functional relationship in which the output can be determined from the input. 
(David Lippman, Business PreCalculus, p1)

\defn{\dt{Function:} \\
	A Function is a rule for a relationship between an input or \underline{\hskip2in} quantity \\ \\
	and an output or \underline{\hskip1.75in} in which each input value \underline{\hskip1.75in} \\ \\
	determines one output value.	We say ``the output is a function of the input''.
}

\ex{Determine the following:
	\begin{enumerate}
		\item Age/Height:
		\begin{enumerate}
			\item Is your height a function of your age? \hrulefill \\
			(\emph{Given your age, can you uniquely determinely your height at that given age?})
			
			\item Is your age a function of your height? \hrulefill \\
			(\emph{Given a specific height, can you uniquely determine your age?}) 
		\end{enumerate}
		
		\item At a coffee shop, the menu consists of items and their prices.
		\begin{enumerate}
			\item Is price a function of the item? \hrulefill \\
			(\emph{Given an item, is their a unique price associated with the given item?})
			
			\item Is the item a function of the price? \hrulefill \\
			(\emph{Given a price, is their a unique item associated with the given price?})
		\end{enumerate}
		
		\item Bank Account:
		\begin{enumerate}
			\item Is your balance a function of your bank account number? \hrulefill \\
			(\emph{If you input a bank account number, does it make sense that the output is your balance?})
			
			\item Is your bank account number a function of your balance? \hrulefill \\
			(\emph{If you input a balance, does it make sense that the output is your bank account number?})
		\end{enumerate}
		
		
	\end{enumerate}
	
}


\pagebreak

\myhead{Function Notation}

To simplify writing out expressions and equations involving functions, a simplified notation is often used.  We also use descriptive variables to help us remember the meaning of the quantities in the problem.

\ex{Rather than write ``height is a function of age'', we could use the descriptive variable $h$ to represent height and we could use the descriptive variable $a$ to represent age.
	
\begin{tabular}{ll}
	``height is a function of age'' & if we name the function $f$ we write \\
	``$h$ is $f$ of $a$'' & or more simply \\
	$h=f(a)$ & we could instead name the function $h$ and write \\
	$h(a)$ & which is read ``$h$ of $a$''
\end{tabular}

\vskip12pt

Words of Advice for Function Notation:
\begin{itemize}
	\item Remember we can use \emph{any} variable to name the function; the notation $h(a)$ shows us that \textbf{$h$ depends on $a$}. The value ``$a$'' \underline{must be} put into the function ``$h$'' to get a result.  
	\item Be careful -- the parentheses indicate that age is input into the function. \\
	(Note: Do not confuse these parentheses with multiplication!)
	
\end{itemize}

}

\defn{\dt{Function Notation:} 
	
	The notation defines a function named $f$.
	\begin{center}
		output = $f$(input)
	\end{center}
	
	This would be read ``output is $f$ of input''.
}


\ex{A function $N = f(y)$ gives the number of police officers, $N$, in a town in year $y$.  \\
	What does $f(2005) = 300$ tell us? \vfill 
}

\pagebreak


\myhead{Evaluating Functions}

Evaluate the following function at the specified values of the independent variable and simplify the results.

\begin{enumerate}
	\item \(f(x) = 4x - 5\) 
	\begin{multicols}{2}
		\begin{enumerate}
			\item
			\(f(1) =\) 
			\vskip0.5in
			
			\item 
			\(f\left( \frac{1}{4} \right) =\)  
			\vskip1in \null
			
			\item
			\(f( - 3) =\) 
			\vskip0.5in
			
			\item
			\(f(x - 1) =\) 
			\vskip1in \null 
		\end{enumerate}
	\end{multicols}
	
	
	\item Let \(f(x) = 4x - 5\) and \(g(x) = 2x + 3\) 
	
	\begin{enumerate}[resume]
		\item
		\((f + g)(2) =\) \vfill 
		\item $(f - g)(3) = $ \vfill 
		\item
		\((f - g)(5a) =\) \vfill 
		\item $(fg)( - 1) = $ \vfill 
	\end{enumerate}
	
\end{enumerate}


%%%% I am cutting this out because we are not longer doing the Limit Def of Derivative.
%
%   \fact{A foundational concept in calculus is the {difference quotient}.  Given a function $f(x)$, where $f$ is defined at $x$ and $x+h$ and $h \neq 0$, the \dt{difference quotient} is  \[\dfrac{f(x + h) - f(x)}{h}\] }
%    
%    Given the function \(f(x) = 7x+2\), find the difference quotient: 
%    
%    $\dfrac{f(3 + h) - f(3)}{h} = $ \vfill 
%\vfill    
%\pagebreak





\pagebreak


\myhead{Additional Practice}  
  
  Evaluate the given function for \(f(x) = x^{2} + 1\) and \(g(x) = x - 4\).
  
  \begin{enumerate}
    \item
      \((f + g)(5) =\)\vfill 
    \item
      \((f - g)(2c) =\) \vfill 
    \item
      \(4 \cdot g( - 3) =\) 
      \vfill 
      

    \item 
      \(3 \cdot f(4) - 2 \cdot g( - 1) =\) 
      \vfill 
      
      \item
      \((fg)( - 2) =\) 
      \vfill 
      
    \item 
     \(\left( \frac{f}{g} \right)(0) =\) 
     \vfill 
     

%      \item 
%      \(\dfrac{f(4) - g(3)}{f(2)} =\) 
%      \vfill 
%      
%    \item 
%      \(\dfrac{g( - 1 + h) - g( - 1)}{h} =\) 
%      \vfill 
%      
%    \item
%     \(\dfrac{f(3 + h) - f(3)}{h}=\) \vfill 
  \end{enumerate}


\end{document}




%%%%%%%%%%%%%%%%%%%%%%%%%%%

\pagebreak

Many situations can be modeled using linear equations and inequalities.  Model each scenario below with a linear equation or inequality.  Then solve the equation or inequality and answer the question.

\begin{enumerate}
	\item Alex purchases a big screen TV, pays 7\% sales tax and is charged \$65 for delivery. If Alex's total cost is \$1668.93, what is the purchase price of the TV?\vfill
	
	\item \emph{Depreciation}. A new car worth \$45,000 is depreciating (decreasing) in value by \$5,200 per year. Find the linear model for the current value, $v$, of the car $y$ years after it was purchased. When will the car be worth \$20,000? \vfill
	
	\item Rita invested a total of $\$60,000$
	in two accounts. One account pays $3.7\%$
	interest annually; the other pays $4.2\%$
	interest annually. At the end of the year, Rita earned a total interest of $\$2,280$. How much money did Rita invest in each account? \vfill
	
\end{enumerate}

\pagebreak

\myhead{Break-even Analysis}

In a business there are two types of costs: fixed costs and variable costs.  

{\bf Fixed costs} are costs that do not change with the number of items produced.  

Examples of fixed costs: \\ \\

{\bf Variable costs} are costs that change with the number of items produced.  

Examples of variable costs: \\ \\

The {\bf total cost} is the sum of the fixed costs and the variable costs.

\hrulefill

\begin{enumerate}[resume]
	\item A publisher for a promising new novel figures fixed costs (overhead, advances, promotion, copyediting, typesetting, and so on) at \$87,000 and variable costs (printing, paper, binding, shipping) at \$4.50 for each book produced. What is the total cost for $x$ books? \vfill
\end{enumerate}



{\bf Revenue} is the amount of money a business receives from selling its products.  This is the number of items sold times the price per item.


The {\bf profit} is the difference between the revenue and the total cost.

If revenue is greater than total cost, then the business makes a profit.  If revenue is less than total cost, then the business loses money.  If revenue is equal to total cost, then the business breaks even.

\begin{enumerate}[resume]
	\item Consider the novel we discussed above. If the book is sold to distributors for \$28 each, how many must be produced and sold for the publisher to break even?
	\vfill
	
	\item A band is planning a concert. The band will receive \$2000 from the venue and \$2 for each ticket sold. The band's expenses are \$2500. How many tickets must be sold for the band to break even? \vfill
\end{enumerate}

\pagebreak

\myhead{Additional Practice}
\begin{enumerate}
	\item Mary paid 8.5\% sales tax and a \$190 title and license fee when she bought a new car for a total of \$28,400. What is the purchase price of the car? \vfill
	\item How many bike computers would a company have to make and sell to break even if the fixed costs are \$36,000, variable costs are \$10.40 per computer, and the computers are sold to retailers for \$15.20 each? \vfill
	\item A company that manufactures and sells x-ray machines has fixed costs of \$1,200,000 and variable costs of \$1,800 per machine. If the machines are sold to hospitals for \$2,400 each, how many machines must be sold for the company to break even? \vfill
\end{enumerate}


